\section{Introduction}
\label{2402:sec:main:introduction}

Recent years have witnessed significant theoretical advancements in understanding the dynamics of training neural networks using gradient descent, to unravel the learning mechanisms of these networks, particularly how they adapt to data and identify pivotal features for predicting the target function. Significant progress has been made over the last few years in the case of two-layer networks, in large part thanks to the so-called mean-field analysis \citep{mei2018mean,chizat2018global,rotskoff2022trainability,sirignano2020mean}). Most of the theoretical efforts, in particular, focused either on {\it one-pass} optimization algorithms, where {\it each iteration involves a new fresh batch of data}, or to the limit of gradient flow in the population loss. For high-dimensional synthetic Gaussian data, and a low dimensional target function (a multi-index model), the class of functions efficiently learned by these \textit{one-pass} methods has been thoroughly analyzed in a series of recent works, and have been shown to be limited by the so-called information exponent \citep{arous2021online} and leap exponent \citep{abbe2022merged,abbe2023sgd} of the target. These analyses have sparked many follow-up theoretical works over the last few months, see, e.g. \citep{damian2022neural,damian2023smoothing,dandi2023how,bietti2023learning,ba2023learning,moniri2023theory,mousavihosseini2023gradient,zweig2023symmetric}. 


However, a common practice in machine learning involves repeatedly traversing the same mini-batch of data. This paper, therefore, aims to go beyond the constraints of single-pass algorithms and to evaluate whether multiple-pass training overcomes these inherent flaws of single-pass methods. We focus on gradient descent, certainly the most straightforward procedure in this family. The theoretical framework we use to prove our main results is based on Dynamical Mean Field Theory (DMFT), which was developed in the statistical physics community \citep{sompolinsky1988chaos} to analyze correlated systems, and recently made rigorous in the context of high dimensional machine-learning problems in \citep{celentano2021high,gerbelot2022rigorous}.

Our findings significantly alter the prevailing narrative in the literature. We demonstrate that gradient descent {\it surpasses} the limitations imposed by the information and leap exponents, achieving a positive correlation with the target function for a much broader class than the staircase functions \citep{abbe2021staircase}, even with minimal (that is, two) repetition of data batches. We characterize the (broad) class of functions efficiently learned in finite time. Among the exceptions are {\it symmetric} functions that remain a challenge due to their extended symmetry-breaking times (a natural feature of physical dynamics \citep{bouchaud1998out}). 

With independent Gaussian datapoints as inputs, both one-pass SGD and gradient flow on population loss result in pre-activations remaining distributed as Gaussian random variables. This Gaussianity underlies the analysis of such settings,  starting from the seminal work of \cite{saad1995line}.
In contrast, upon re-using batches, the pre-activations develop non-Gaussian components correlated with the targets.  This non-Gaussianity is a crucial aspect of the stark contrast in the learning of directions compared to the one-pass setting. While we establish our results for discrete steps with extensive batches of size $\cO(d)$ where $d$ denotes the input dimension, we expect our conclusions about the learning of new directions to also hold while performing gradient flow on the empirical loss, since such a setup would lead to the development of similar non-Gaussian components in the pre-activations, in contrast to gradient flow on the population loss where the pre-activations remain Gaussian.

Our results demonstrate that contrary to the common wisdom ``the more data the better", gradient descent on the same batch can surpass one-pass SGD on different batches, even when one-pass SGD utilizes a larger number of datapoints. More generally, we believe that our analysis provides insights into incremental learning of features in the presence of correlations between datapoints across batches, which is typical in most high dimensional datasets having a small number of ``latent" factors. Our conclusions follow from a rigorous mathematical proof rooted in DMFT, which we also use to provide an analytic description of the dynamic processes of low-dimensional weight projections. This analysis has interest on its own.
\looseness=-1
